\RequirePackage{fix-cm}

%\documentclass{svjour3}                     % onecolumn (standard format)
%\documentclass[smallcondensed]{svjour3}     % onecolumn (ditto)
\documentclass[smallextended]{svjour3}       % onecolumn (second format)
%\documentclass[twocolumn]{svjour3}          % twocolumn
%
\smartqed  % flush right qed marks, e.g. at end of proof
%
\usepackage{graphicx}
\usepackage{amsmath}
\usepackage{amssymb}
\usepackage{booktabs}
\usepackage{algorithm}
\usepackage{algorithmic}
\usepackage{multirow}
\usepackage{tabularx}


\begin{document}

\title{Modeling and Optimizing Mixed AI-Human Software Development Teams:
       A Stochastic Simulation Framework}
%\titlerunning{Optimizing Mixed AI-Human Development Teams}

\author{Steffen Herbold}

%\authorrunning{Herbold}

\institute{Steffen Herbold, 0000-0001-9765-2803 \at
University of Passau, Passau, Germany.
\email{steffen.herbold@uni-passau.de}
}

\date{Received: date / Accepted: date}
% The correct dates will be entered by the editor


\maketitle

\begin{abstract}
% --- ABSTRACT (write last) ---
% Structure: problem -- gap -- approach -- model highlights -- key findings -- implications
%
% - Problem: AI agents are increasingly deployed alongside human developers, yet no
%   principled quantitative framework guides how many agents of each type to use.
% - Gap: existing SE process models treat team headcount as a scalar and ignore
%   role-specific quality degradation, CI/CD latency effects, and multi-objective trade-offs.
% - Approach: we present a formal 14-dimensional discrete-time stochastic model of
%   mixed AI-human sprint dynamics, with 11 transition steps, six proved stability
%   conditions, and Monte Carlo + evolutionary search over agent composition vectors.
% - Model highlights: three Pareto objectives (V = throughput, Q = escaped bugs per
%   feature, C_F = cost per feature); 13 agent roles + 2 human roles; four degradation
%   functions linking technical debt and review coverage to quality and productivity.
% - Key findings: summarize after experiments are run (Pareto frontier shape,
%   dominant configurations, sensitivity insights).
% - Implications: practical guidance on minimum Code Review / Testing allocation,
%   Refactoring-to-Implementation ratio, and span-of-control limits.
%
% Target length: 150--200 words.
TODO: write abstract.
\end{abstract}

\keywords{AI agents \and software process simulation \and multi-objective optimization
\and technical debt \and Pareto frontier \and team composition}


% ============================================================
\section{Introduction}
\label{sec:introduction}
% ============================================================
%
% CONTENT OUTLINE
%
% [Motivation / opening hook]
% - Software development teams increasingly mix human developers with autonomous AI
%   agents (coding assistants, automated reviewers, test generators, etc.).
% - The central practical question: how many agents of each type should be deployed,
%   and how many human developers are needed to supervise them?
% - No principled quantitative framework currently exists to answer this question;
%   practitioners rely on intuition or ad-hoc experiments.
%
% [Core tension]
% - More agents raise throughput but also raise coordination cost, review load,
%   technical debt pressure, and CI/CD latency -- objectives conflict and cannot
%   be optimized in isolation.
% - Human developers are a bottleneck resource: span-of-control M (agents per human)
%   is bounded by human capacity, review quality degradation, and escalation load.
%
% [Research questions]
% - RQ1: What agent compositions are Pareto-optimal on (V, Q, C_F) under realistic
%        sprint dynamics?
% - RQ2: How do degradation effects (technical debt, review quality decay,
%        escalation overload) shift the Pareto frontier relative to a simplified
%        no-degradation baseline?
% - RQ3: How sensitive are optimal configurations to workload scenario
%        (default / high-load / low-budget)?
%
% [Contributions -- four bullet points]
% 1. Formal 14-dimensional state-space model of AI-human sprint dynamics with
%    11 ordered transition steps and proved convergence properties.
% 2. Six stability/feasibility theorems establishing conditions for bounded
%    feature backlog, bug queue, human capacity, and technical debt.
% 3. Monte Carlo + evolutionary Pareto search over the agent composition space,
%    with a deterministic feasibility pre-filter to prune infeasible configurations.
% 4. Empirical results across three parameterised scenarios identifying which
%    agent roles drive each Pareto objective and quantifying degradation effects.
%
% [Optional teaser figure]
% - Figure 1: schematic of a mixed team (human developers supervising AI agents)
%   and a sketch of the resulting tri-objective Pareto frontier.
%
% [Paper structure -- one sentence per section]



% ============================================================
\section{Background and Related Work}
\label{sec:related}
% ============================================================

% ---- 2.1 ----
\subsection{Software Team and Process Simulation}
\label{sec:related:process}
%
% - Situate within software process simulation (SPS): Abdel-Hamid \& Madnick
%   system-dynamics model, COCOMO-derivative planning tools, agent-based SE
%   process models (e.g., Concepcion \& Hilburn).
% - Prior models treat headcount as a scalar; this work parameterises role
%   composition as a 13-dimensional integer vector.
% - Discrete-time sprint framing is closer to Scrum simulation literature
%   (Cao et al.) than to continuous-time system dynamics.
% - Limitations of prior models for AI-agent settings: uniform human productivity
%   assumption, no CI/CD latency buffer, no role-specific degradation functions.
% - Brooks' Law formulation: H_comm(N) = c_comm * (N-1) communication overhead
%   is incorporated explicitly in the human capacity model.

% ---- 2.2 ----
\subsection{AI Agents in Software Engineering}
\label{sec:related:agents}
%
% - Empirical evidence on AI-assisted coding: GitHub Copilot productivity studies,
%   LLM-based code review (CodeReviewer, etc.), automated bug fixing (SWE-bench,
%   Devin, SWE-agent), automated test generation (CodaMosa, ChatUniTest).
% - Distinguish tool-assisted human work from autonomous agent roles as modelled
%   here: agents execute full tasks independently, not merely assist.
% - Gap: individual-agent productivity studies do not model systemic effects on
%   backlog stability, bug escape rates, or cost per feature in mixed teams.
% - Emerging multi-agent SE frameworks (ChatDev, MetaGPT, OpenHands) as
%   direct industrial motivation for the team-composition question.
% - Cost structure reality: LLM API token pricing maps naturally to per-agent
%   per-sprint cost model used here.

% ---- 2.3 ----
\subsection{Multi-Objective Optimisation in Software Engineering}
\label{sec:related:moea}
%
% - SBSE (search-based SE) for team and project planning: staffing optimisation,
%   sprint planning, release planning under uncertainty.
% - Pareto-based formulations in SE: NSGA-II/III applied to test prioritisation,
%   configuration optimisation, architecture selection.
% - Key distinction: our feasibility pre-filter (six stability conditions at prior
%   means) dramatically prunes the search space before any simulation, a technique
%   not present in prior SBSE team-composition work.
% - Motivate random + evolutionary hybrid search as appropriate for discrete,
%   high-dimensional, partially infeasible composition spaces.
% - Hypervolume indicator as convergence metric (Zitzler \& Thiele 1998).


% ============================================================
\section{The Formal Model}
\label{sec:model}
% ============================================================

% ---- 3.1 ----
\subsection{Sprint and Release Cycle}
\label{sec:model:cycle}

The model treats software development as a discrete-time stochastic dynamical
system. Time is indexed by sprints $t \in \{0, 1, \ldots, T-1\}$; each sprint
represents one fixed-duration iteration. The
simulation horizon $T$ and the team configuration
$(N, M, \mathbf{A}_\mathrm{dev}, \mathbf{A}_\mathrm{qa}, R, F_\mathrm{freeze})$
are fixed across all $K$ independent Monte Carlo replications of a single run.
Within each replication, \emph{scenario-level} parameters (throughput rates,
bug injection rates, detection probabilities, cost coefficients) are sampled
once at $t = 0$ and held constant for all $T$ sprints; \emph{step-level}
randomness (Poisson feature arrivals, Binomial detection draws) is resampled at
every sprint. This two-level sampling structure separates between-team variation
(captured by different configurations) from within-team uncertainty (captured
across $K$ replications).

\paragraph{Within-sprint dynamics.}
At each sprint the system executes eleven ordered steps covering feature
arrivals and implementation, bug injection and Stage-A detection, Stage-B
latency promotion and detection, architecture review and escalation queue
updates, bug fixing, regression injection, technical debt evolution, test-suite
and documentation growth, CI/CD latency recomputation, and release processing
(Section~\ref{sec:model:transition}). All quantities within a sprint are
computed as aggregate counts; within-sprint ordering among agents is not
modelled. Agents of the same type are
statistically identical within a replication, so the composition count
$a_k$ is the only role-level degree of freedom. Each sprint selects the active
agent composition based on its position in the release cycle: the
\emph{development composition} $\mathbf{A}_\mathrm{dev}$ is used during the
first $R - F_\mathrm{freeze}$ sprints of each cycle; the \emph{QA composition}
$\mathbf{A}_\mathrm{qa}$ is used during the remaining $F_\mathrm{freeze}$
sprints (see below).

\paragraph{Two-stage detection pipeline.}
Bug detection operates on two timescales that are central to the model's
quality dynamics. \emph{Stage A} (Code Review and Security Audit) is
synchronous: it acts within the same sprint a bug is introduced, because
reviewers inspect the code change directly. Bugs surviving Stage A enter a
per-severity FIFO latency buffer and are not testable until
$\lceil\tau_\mathrm{CI}(t)\rceil$ sprints have elapsed, at which point
\emph{Stage B} (Unit, Integration, and System/Acceptance tests) processes them.
Regression bugs introduced by fixing or refactoring skip Stage A and enter the
buffer immediately. This asymmetry means that increasing code-review
coverage and increasing test coverage are \emph{not} interchangeable: Stage A
prevents bugs from entering the hidden stock at all, whereas Stage B reduces the
hidden stock with a delay that grows with codebase and test-suite size.

\paragraph{Release events and feature freeze.}
A release occurs exactly at sprint indices $R, 2R, 3R, \ldots$ where $R \geq 1$
is the fixed release cadence. At each release all bugs remaining in the
latency buffer (injected within the last $\lceil\tau_\mathrm{CI}\rceil$ sprints
and therefore untested) and all Stage-B survivors (bugs that exited the latency
buffer but evaded all three test layers) are counted as production escapes and
added to $B_\mathrm{escaped}$. 

Similar to real-world projects, a feature freeze is implemented in the model to
stabilise the codebase before release. Additionally, this mechanism ensures that
bugs injected near the end of a release cycle have time to propagate through the
Stage-B pipeline before the release. The feature freeze lasts for 
$F_\mathrm{freeze} \in [1, R-1]$ sprints before a release.
Sprint $t$ is in the freeze window iff $(t \bmod R) \geq R - F_\mathrm{freeze}$.
During these sprints the team switches from $\mathbf{A}_\mathrm{dev}$ to
$\mathbf{A}_\mathrm{qa}$: the QA composition sets $a_1 = a_2 = a_3 = 0$
(no new requirements, designs, or features) and redeploys the full $N \times M$
agent budget to bug-fixing, testing, refactoring, and documentation roles.
Human direct-coding capacity is also suspended during the freeze so that the
residual human capacity is fully available for review and escalation work.

% TODO: Figure -- sprint cycle and release timeline diagram

% ---- 3.2 ----
\subsection{Roles}
\label{sec:model:roles}

\subsubsection{Human roles}
\label{sec:model:roles:human}

Two human roles are included in the model. The \textbf{Product Owner} is a
singleton who provides raw feature requirements, prioritises the backlog,
defines acceptance criteria, and authorises releases. It is modelled as a
fixed-overhead resource carrying an annual cost of $c_\mathrm{PO}$; new
features arrive to the backlog as a Poisson process with rate $\lambda_F$
features per sprint, independently of team composition. The \textbf{Human
Developer} is the central scalable human role. Each of the $N$ developers
supervises $M$ agents, reviews and approves architecture proposals, resolves
escalations that agents cannot handle autonomously, and contributes direct code
to tasks that exceed agent capability. Each developer carries an annual
fully-loaded salary $c_\mathrm{salary}$ and a per-sprint time budget of
$H_\mathrm{capacity}$ hours. This budget is consumed in strict priority order by
(i)~peer communication overhead
$H_\mathrm{comm}(N) = c_\mathrm{comm}(N-1)$,
capturing Brooks' coordination cost~\cite{Brooks1975}, (ii)~agent monitoring,
(iii)~architecture review queue service, (iv)~escalation queue service, and
finally (v)~residual direct coding. As the span of control $M$ increases,
monitoring overhead grows linearly and crowds out productive capacity, creating
the fundamental throughput--quality--cost trade-off that the simulation is
designed to characterise. The full capacity allocation model is given in
Section~\ref{sec:model:transition}.

\subsubsection{AI agent roles}
\label{sec:model:roles:agent}

Thirteen specialist agent types are modelled, indexed $k = 1,\ldots,13$.
Each agent of type $k$ executes tasks autonomously at a per-sprint throughput
rate $r_k$ drawn from a log-normal distribution, reflecting the empirically
observed right-skew of software task durations~\cite{TODO:cite}. Routine
outputs require no human approval; agents escalate to their supervising \textit{Human
Developer} only when a task exceeds their capability, at a per-agent escalation
rate $\rho_\mathrm{esc}$. Each deployed agent of type $k$ incurs a per-sprint
API token cost $c_{\mathrm{token},k}$, reflecting current LLM pricing;
agent costs therefore scale with deployment count whereas human costs are
fixed salaries, making the human-to-agent ratio a primary lever for the cost
objective $C_F$. Table~\ref{tab:roles} summarises each role's primary function,
principal effect on the simulation state, and position in the detection
pipeline.

The thirteen roles fall into four functional groups. \emph{Production roles}
(k = 1--5) drive or shape the feature pipeline: the Requirements Refinement
Agent refines specifications (affecting the specification-error amplifier
$\varepsilon_\mathrm{spec}$), the Architecture Agent generates design proposals
that feed the human review queue, the Feature Implementation Agent is the
primary productivity driver, the Bug Fixing Agent drains the known-bug queue,
and the Refactoring Agent is the only role that actively \emph{reduces}
technical debt $D(t)$. Without adequate refactoring capacity, debt converges
to a positive equilibrium that permanently slows implementation and amplifies
bug injection (Section~\ref{sec:stability:debt}). \emph{Stage-A detection
roles} (k = 6--7) inspect every code change synchronously within the sprint it
is produced: the Code Review Agent detects non-security defects at effective
rate $d_\mathrm{CR,eff}$, and the Security Audit Agent detects security-class
defects at rate $d_\mathrm{SA}$; bugs surviving Stage A enter the CI/CD
latency buffer and become detectable only after $\lceil\tau_\mathrm{CI}(t)\rceil$
sprints. \emph{Stage-B detection roles} (k = 8--10) build up automated test
suites whose coverage grows with $N_\mathrm{UT}$, $N_\mathrm{IT}$,
$N_\mathrm{ST}$ respectively, providing delayed but increasingly thorough
detection as the project matures. \emph{Support roles} (k = 11--13) affect
operational overhead and cost: the DevOps Agent reduces the base CI/CD latency
$\tau_\mathrm{CI,base}$; the Developer Documentation Agent grows documentation
coverage $\mathrm{doc\_cov}(t)$, which in turn reduces the per-agent monitoring
overhead that consumes human capacity; and the User Documentation Agent
contributes a per-release cost $C_\mathrm{user\_doc}$ that scales with the
features shipped in each release.

One modelling assumption deserves explicit note. Code Review throughput is
treated as \emph{unbounded}: all changed
code is reviewed regardless of $a_6$, but effective detection quality
$d_\mathrm{CR,eff}$ degrades with change volume and with architecture-review
queue pressure when $a_6$ is small. Each deployed Code Review Agent still
accrues its full token cost $c_{\mathrm{token},6}$, so over-allocating agents
to this role is penalised by the cost objective; Pareto optimality therefore
imposes an implicit upper bound on $a_6$.

\begin{table}[t]
\caption{The thirteen AI agent roles. \emph{Det.\ stage}: A\,=\,synchronous
  same-sprint detection; B\,=\,delayed detection via CI/CD latency buffer;
  ---\,=\,no direct detection function.}
\label{tab:roles}
\footnotesize
\begin{tabularx}{\linewidth}{cp{2.6cm}Xp{3.0cm}c}
\toprule
$k$ & Role & Primary Function & Principal State Effect & Det. \\
\midrule
1 & Requirements Refinement
  & Transforms ambiguous requirements into unambiguous specifications
  & Modulates spec-error amplifier $\varepsilon_\mathrm{spec}$ on bug injection
  & --- \\[2pt]
2 & Architecture \& Design
  & Produces component designs and interface contracts for human approval
  & Drives $q_\mathrm{arch}$; proposals enter review queue $Q_\mathrm{rev}$
  & --- \\[2pt]
3 & Feature Implementation
  & Implements features per approved designs
  & Drains $F_\mathrm{backlog}$; increments $F_\mathrm{completed}$; injects bugs; accumulates debt via $\delta_\mathrm{impl}$
  & --- \\[2pt]
4 & Bug Fixing
  & Diagnoses and repairs defects reported by the detection pipeline
  & Drains $B_\mathrm{known}$; regressions (prob.\ $p_\mathrm{regress}$) bypass Stage~A
  & --- \\[2pt]
5 & Refactoring
  & Restructures code to reduce complexity without changing behaviour
  & Decreases $D(t)$ at rate $\delta_\mathrm{refactor}$; regression risk $p_\mathrm{regress\_refactor}$
  & --- \\[2pt]
6 & Code Review
  & Reviews all code changes for correctness, style, and design
  & Synchronously detects fraction $d_\mathrm{CR,eff}(t)$ of non-security bugs
  & A \\[2pt]
7 & Security Audit
  & Identifies security-class vulnerabilities in code and design
  & Synchronously detects fraction $d_\mathrm{SA}$ of security bugs
  & A \\[2pt]
8 & Unit Test Development
  & Writes automated unit tests; grows coverage over time
  & Accumulates $N_\mathrm{UT}$; Stage-B detection $d_\mathrm{UT,eff}$
  & B \\[2pt]
9 & Integration Test Development
  & Tests cross-component interactions and interface contracts
  & Accumulates $N_\mathrm{IT}$; Stage-B detection $d_\mathrm{IT,eff}$
  & B \\[2pt]
10 & System \& Acceptance Testing
  & Validates end-to-end behaviour against acceptance criteria
  & Accumulates $N_\mathrm{ST}$; Stage-B detection $d_\mathrm{ST,eff}$; final release gate
  & B \\[2pt]
11 & DevOps / CI-CD
  & Builds and maintains the build, test, and deployment pipeline
  & Reduces $\tau_\mathrm{CI,base}$; incurs $C_\mathrm{DevOps,setup}$ and $C_\mathrm{DevOps,maint}$
  & --- \\[2pt]
12 & Developer Documentation
  & Produces API references, architecture guides, and contributor docs
  & Grows $\mathrm{doc\_cov}(t)$; reduces per-agent monitoring overhead via $\Gamma_\mathrm{doc}$
  & --- \\[2pt]
13 & User Documentation
  & Produces user manuals, tutorials, and release notes
  & Contributes $C_\mathrm{user\_doc}$ per release, proportional to features shipped
  & --- \\
\bottomrule
\end{tabularx}
\end{table}

% ---- 3.3 ----
\subsection{Project and team structure}
\label{sec:model:team}

A the project and team structure is fully specified by six control variables: the human developer count
$N \geq 1$, the number of agents each developer can control $M \geq 1$, the number of sprints per release $R \geq 1$, the length of the feature freeze period $F_\mathrm{freeze} \in
[1, R]$, and two 13-element agent composition vectors $\mathbf{A}_\mathrm{dev}$ and
$\mathbf{A}_\mathrm{qa}$ described below.

$\mathbf{A}_\mathrm{dev} = (a_{dev,1},\ldots,a_{dev,13}) \in \mathbb{N}^{13}$ is the
\emph{development-phase composition}: the count of agents of each type deployed
during the $R - F_\mathrm{freeze}$ active development sprints of each release
cycle. $\mathbf{A}_\mathrm{qa} = (a_{qa,1},\ldots,a_{qa,13}) \in \mathbb{N}^{13}$ is
the \emph{QA-phase composition}: the agent allocation during the
$F_\mathrm{freeze}$ stabilisation sprints. Both vectors satisfy the same
\emph{budget constraint}
\begin{equation}
  \sum_{k=1}^{13} a_{\cdot,k} = N \times M,
  \label{eq:budget}
\end{equation}
ensuring the total agent headcount is fixed throughout the release cycle;
redeployment from development to QA roles incurs no extra cost.
$\mathbf{A}_\mathrm{dev}$ must satisfy $a_{dev,3} \geq 1$ (at least one Feature
Implementation Agent in the development phase). $\mathbf{A}_\mathrm{qa}$ is
required to have $a_{qa,1} = a_{qa,2} = a_{qa,3} = 0$, to ensure that no new features are added during a feature freeze and redeploy all agents to QA roles. All $a_{\cdot,k}$ agents of the
same type are assumed to be statistically identical.



% ---- 3.4 ----
\subsection{State Space}
\label{sec:model:state}

The system state at the start of sprint $t$ is a 14-component tuple $S(t)$,
updated at the end of each sprint by the eleven-step transition function
described in Section~\ref{sec:model:transition}. The components fall into five
clusters summarised in Table~\ref{tab:state}.

\subsubsection{Feature pipeline.}
$F_\mathrm{backlog}(t) \in \mathbb{N}_{\geq 0}$ counts features awaiting
implementation at the start of sprint $t$; it grows by $\mathrm{Poisson}(\lambda_F)$
arrivals each sprint and is drained by the \textit{Feature Implementation Agents} and the
direct coding capacity of \textit{Human Developers}. $F_\mathrm{completed}(t) \in \mathbb{N}_{\geq 0}$
is the cumulative count of fully implemented features; it is non-decreasing and
appears as the normalising denominator in all three Pareto objectives.

\subsubsection{Bug stocks.}
Bugs are tracked at three pipeline stages --- hidden, known, and escaped --- and
at each stage are split by severity: minor (fraction $1-\theta$ of injected
bugs) and major (fraction $\theta$). The severity split is essential because
minor and major bugs have different fix rates, fix costs, and production impact.

$B_\mathrm{hid,min}(t)$ and $B_\mathrm{hid,maj}(t)$ count bugs that exist in
the codebase but have not yet been detected by any review or test layer; they
equal the total size of the CI/CD latency buffers described below.
$B_\mathrm{kno,min}(t)$ and $B_\mathrm{kno,maj}(t)$ count bugs that have been
detected and are queued for repair by the Bug Fixing Agent. The separation
between hidden and known is load-bearing: it ensures that fixing capacity is
consumed only by detected bugs, while hidden bugs silently degrade quality until
they are either caught or escape to production. $B_\mathrm{escaped}(t) \in
\mathbb{N}_{\geq 0}$ is the cumulative count of bugs that reached production
without detection; it grows only at release events (every $R$ sprints).

\subsubsection{Technical debt and human capacity queues.}
$D(t) \in [0, 1]$ is a normalised measure of technical debt: $D = 0$ represents
a debt-free codebase and $D = 1$ represents maximum structural degradation. It
is the principal coupling variable in the model: through the debt amplifier
$\phi_D(D)$ it multiplicatively increases the bug injection rate per feature,
and through the debt slowdown $\psi_D(D)$ it multiplicatively reduces
implementation and fix throughput (Section~\ref{sec:model:degradation}). The
\textit{Refactoring Agent} is the only role that reduces $D(t)$.
$Q_\mathrm{rev}(t) \in \mathbb{N}_{\geq 0}$ and $Q_\mathrm{esc}(t) \in
\mathbb{N}_{\geq 0}$ are the lengths of the architecture proposal review backlog
and the agent escalation queue respectively; both are bottleneck state variables
that consume \textit{Human Developer} capacity and trigger review- and escalation-quality
degradation when they grow large (Section~\ref{sec:model:degradation}).

\subsubsection{Test infrastructure and documentation.}
$N_\mathrm{UT}(t)$, $N_\mathrm{IT}(t)$, and $N_\mathrm{ST}(t)$ are the
cumulative numbers of unit, integration, and system/acceptance tests written up
to sprint $t$; they grow monotonically as the respective testing agents produce
new test cases. Test \emph{coverage} is derived from these counts relative to
$F_\mathrm{completed}(t)$ (see derived quantities below) and determines the
effective Stage-B detection probability. Storing raw counts rather than coverage
fractions avoids a circular dependency in the sprint transition, since both
numerator and denominator change within the same sprint.
$\mathrm{doc\_cov}(t) \in [0, 1]$ is the fraction of the codebase covered by
developer documentation; it erodes each sprint by a factor proportional to new
features that ship without immediate documentation, and grows when the Developer
Documentation Agent is active. Documentation coverage reduces per-agent
monitoring overhead through a reduction factor $\Gamma_\mathrm{doc}$, linking
investment in documentation to human capacity savings.

\begin{table}[t]
\caption{State vector components with domains and the sprint-transition
  steps (cf.\ Section~\ref{sec:model:transition}) that write each variable.}
\label{tab:state}
\footnotesize
\begin{tabularx}{\linewidth}{p{2.3cm}Xp{1.5cm}p{1.4cm}}
\toprule
Symbol & Description & Domain & Modified (step) \\
\midrule
$F_\mathrm{backlog}(t)$     & Features awaiting implementation                      & $\mathbb{N}_{\geq 0}$  & 1 \\[2pt]
$F_\mathrm{completed}(t)$   & Cumulative features shipped; non-decreasing            & $\mathbb{N}_{\geq 0}$  & 1 \\[2pt]
$B_\mathrm{hid,min}(t)$     & Hidden minor bugs (latency-buffer total)               & $\mathbb{N}_{\geq 0}$  & 2, 3, 8 \\[2pt]
$B_\mathrm{hid,maj}(t)$     & Hidden major bugs (latency-buffer total)               & $\mathbb{N}_{\geq 0}$  & 2, 3, 8 \\[2pt]
$B_\mathrm{kno,min}(t)$     & Known minor bugs queued for repair                     & $\mathbb{N}_{\geq 0}$  & 3, 6 \\[2pt]
$B_\mathrm{kno,maj}(t)$     & Known major bugs queued for repair                     & $\mathbb{N}_{\geq 0}$  & 3, 6 \\[2pt]
$B_\mathrm{escaped}(t)$     & Cumulative bugs escaped to production; non-decreasing  & $\mathbb{N}_{\geq 0}$  & 11 \\[2pt]
$D(t)$                       & Technical debt level                                   & $[0, 1]$               & 9 \\[2pt]
$Q_\mathrm{rev}(t)$         & Architecture review queue length                       & $\mathbb{N}_{\geq 0}$  & 4 \\[2pt]
$Q_\mathrm{esc}(t)$         & Escalation queue length                                & $\mathbb{N}_{\geq 0}$  & 5 \\[2pt]
$N_\mathrm{UT}(t)$          & Cumulative unit tests written                          & $\mathbb{R}_{\geq 0}$  & 10 \\[2pt]
$N_\mathrm{IT}(t)$          & Cumulative integration tests written                   & $\mathbb{R}_{\geq 0}$  & 10 \\[2pt]
$N_\mathrm{ST}(t)$          & Cumulative system/acceptance tests written             & $\mathbb{R}_{\geq 0}$  & 10 \\[2pt]
$\mathrm{doc\_cov}(t)$      & Developer documentation coverage                      & $[0, 1]$               & 10 \\
\bottomrule
\end{tabularx}
\end{table}

\subsubsection{Latency buffers.}
Two auxiliary FIFO structures, $\mathrm{buf\_min}$ and $\mathrm{buf\_maj}$, are
maintained alongside $S(t)$ but are not part of it. Each maps an injection
sprint $t_\mathrm{inj}$ to a pair $\bigl(\text{count},\,
\lceil\tau_\mathrm{CI}(t_\mathrm{inj})\rceil\bigr)$, where the second element
is the CI/CD latency snapshot taken at the time of injection. A cohort injected
at sprint $t_\mathrm{inj}$ first becomes detectable by the Stage-B test layers
at any sprint $t$ satisfying $t \geq t_\mathrm{inj} +
\lceil\tau_\mathrm{CI}(t_\mathrm{inj})\rceil$. Storing the \emph{injection-time}
latency per cohort is important: because $\tau_\mathrm{CI}$ grows over the
project lifetime (see below), applying the current latency to old cohorts would
incorrectly delay detection of bugs that were injected when the pipeline was
faster. At each release event all cohorts remaining in the buffers --- too
recently injected to have cleared the pipeline --- are drained directly into
$B_\mathrm{escaped}$. The scalar state components $B_\mathrm{hid,min}(t)$ and
$B_\mathrm{hid,maj}(t)$ are maintained as running totals equal to the respective
buffer sums, avoiding a full buffer traversal at each sprint.

\subsubsection{Derived quantities.}
Three quantities are recomputed from $S(t)$ and the scenario parameters at
each sprint rather than stored as state. The \emph{test coverage levels}
\begin{equation}
  \mathrm{cov}_k(t) = \min\!\left(c_{k,\mathrm{ceil}},\;
    \frac{N_k(t)}{\max\!\bigl(1,\,F_\mathrm{completed}(t)\bigr)}\right),
  \quad k \in \{\mathrm{UT},\,\mathrm{IT},\,\mathrm{ST}\},
  \label{eq:coverage}
\end{equation}
express the ratio of tests written to features shipped, capped at a ceiling
$c_{k,\mathrm{ceil}} \leq 1$ that reflects the practical limits of automated
coverage; they govern the effective Stage-B detection probability
$d_k \times \mathrm{cov}_k(t)$ in Step~3 of the transition. The
\emph{weighted test count}
\begin{equation}
  n_\mathrm{tests}(t) =
    \beta_\mathrm{UT}\,N_\mathrm{UT}(t)
    + \beta_\mathrm{IT}\,N_\mathrm{IT}(t)
    + \beta_\mathrm{ST}\,N_\mathrm{ST}(t),
  \label{eq:ntests}
\end{equation}
with execution-time weights $\beta_\mathrm{UT} \leq \beta_\mathrm{IT} \leq
\beta_\mathrm{ST}$, drives the \emph{CI/CD pipeline latency}
\begin{equation}
  \tau_\mathrm{CI}(t) = \min\!\left(T,\;\;
    \tau_\mathrm{CI,base}
    \times \bigl(1 + \alpha_{F,\mathrm{CI}}\,F_\mathrm{completed}(t)\bigr)
    \times \bigl(1 + \alpha_{\mathrm{test},\mathrm{CI}}\,n_\mathrm{tests}(t)\bigr)%
  \right),
  \label{eq:tau_ci}
\end{equation}
where $\tau_\mathrm{CI,base}$ is the base pipeline latency, and the two
multiplicative factors model the empirically observed growth of CI execution
time with codebase and test-suite size. The cap at $T$ prevents detection from
being delayed beyond the simulation horizon. 

\subsubsection{Initial conditions.}
At $t = 0$ all components of $S(0)$ are zero except $D(0) \sim
\mathrm{Beta}(\alpha_{D_0}, \beta_{D_0})$, sampled once per replication to
represent inherited technical debt. Setting $\alpha_{D_0} \gg \beta_{D_0}$
models a greenfield project with minimal debt; reversing the ratio models a
legacy codebase. An optional warm-start parameter $R_\mathrm{warmup} \in
\mathbb{N}_{\geq 0}$ initialises the feature backlog as $F_\mathrm{backlog}(0)
\sim \mathrm{Poisson}(\lambda_F \times R_\mathrm{warmup})$, representing a team
whose backlog was already accumulating before the simulation window opens.

% ---- 3.5 ----
\subsection{Sprint Transition Function}
\label{sec:model:transition}
%
% Ordered 11-step transition T: S(t) -> S(t+1). Order is not arbitrary:
% latency promotion (step 3) precedes detection (step 8); human capacity
% allocation (step 7) follows agent work (steps 4-6).
%
% Step 1 -- Feature arrival and implementation (Sec. 4.1 of framework):
%   - New features: Poisson(lambda_F) arrivals added to F_backlog.
%   - Implementation blocked if Q_review > threshold.
%   - Throughput: r_F * psi_D(D) * a_3 + H_impl(t) features completed.
%   - Human implementation capacity H_impl(t) = residual human capacity
%     after monitoring/review/escalation (see Step 7).
%
% Step 2 -- Bug injection (Sec. 4.2):
%   - lambda_impl_eff(t) = lambda_impl * phi_D(D) * (1 + alpha_spec*eps_spec)
%     * (2 - q_arch_eff) bugs injected per feature completed.
%   - Severity split: fraction theta -> major; (1-theta) -> minor.
%   - Stage-A detection (Code Review + Security Audit) applied immediately;
%     surviving bugs enter latency buffer buf[t].
%
% Step 3 -- Stage-B detection / latency promotion (Sec. 4.3, 6):
%   - Bugs in buf[t_inj] become detectable at t >= t_inj + ceil(tau_CI(t_inj)).
%   - Detection applied sequentially: Unit tests, Integration tests, System tests.
%   - Coverage-weighted effective detection: d_UT_eff = d_UT * min(cov_UT, cov_ceil_UT).
%   - Detected bugs move from hidden to B_known.
%
% Step 4 -- Architecture review queue update (Sec. 4.4):
%   - q_arch_eff(t) degrades when eta_rev(Q_review) < 1.
%   - Poor architecture quality amplifies bug injection in next sprint via Step 2.
%
% Step 5 -- Escalation queue update (Sec. 4.5):
%   - Failed escalations re-queue: delta_Q_esc = (1 - eta_esc(Q_escalation)) * Q_escalation.
%
% Step 6 -- Bug fixing (Sec. 4.6):
%   - Fixes: min(r_fix * psi_D(D) * a_4, B_known) bugs resolved.
%   - Separate rates for minor (r_fix_min) and major (r_fix_maj) bugs.
%
% Step 7 -- Human capacity allocation (Sec. 3 of framework):
%   Priority-order budget per human per sprint:
%   1. Communication overhead: H_comm(N) = c_comm * (N-1)
%   2. Agent monitoring: M * H_monitoring_overhead * (1 - Gamma_doc * doc_cov)
%   3. Architecture review queue service
%   4. Escalation queue service
%   5. Residual direct coding: H_impl(t)
%
% Step 8 -- Regression injection (Sec. 4.7):
%   - Each fix injects new bug with prob p_regress (enters latency buffer directly,
%     skipping Stage A per assumption A6).
%   - Each refactoring unit injects with prob p_regress_refactor.
%
% Step 9 -- Technical debt update (Sec. 4.8):
%   D(t+1) = clip(D(t) + delta_impl*impl(t) + delta_fix*fixes(t)
%            - delta_refactor * a_5 * r_refactor, 0, 1)
%
% Step 10 -- Test suite and documentation update (Sec. 4.9):
%   - N_UT(t+1) = N_UT(t) + a_8 * r_UT; coverage cov_UT grows with N_UT.
%   - doc_cov(t+1) erodes by factor zeta_doc then grows with a_{12} * r_doc.
%
% Step 11 -- CI/CD latency update (Sec. 4.10, 4.11):
%   - tau_CI(t+1) = tau_CI_base * (1 + alpha_F_CI*F_completed)
%                                * (1 + alpha_test_CI*weighted_tests)
%   - Release event (every R sprints): bugs remaining in latency buffer escape
%     to production; B_escaped updated; user documentation cost incurred.
%
% Figure 3: flowchart of the 11-step transition annotated with state variables
%   modified at each step.

% ---- 3.6 ----
\subsection{Degradation Functions}
\label{sec:model:degradation}
%
% Four functions governing how performance decays under stress:
%
% phi_D(D) = 1 + (phi_max - 1) * D^kappa
%   -- debt AMPLIFIER on bug injection rate (strictly increasing in D).
%   -- phi_max > 1; kappa controls curvature (default kappa=1 => linear).
%   -- Motivation: accumulated debt causes more regressions per change.
%
% psi_D(D) = 1 - (1 - psi_min) * D^gamma
%   -- debt SLOWDOWN on throughput (strictly decreasing; psi_D in [psi_min, 1]).
%   -- gamma controls how quickly slowdown sets in.
%   -- Motivation: navigating complex/messy code takes longer.
%
% eta_rev(Q_review): piecewise linear review quality decay
%   -- eta_rev = 1 when Q_review <= Q_review_threshold (queue manageable).
%   -- Decays linearly to eta_min_rev as queue grows beyond threshold.
%   -- Motivation: reviewers rush when backlog is large; miss more defects.
%
% eta_esc(Q_escalation): piecewise linear escalation quality decay
%   -- Analogous form to eta_rev.
%   -- Poor resolutions cause re-escalation (feeds back into Q_escalation).
%
% - Discuss why all four are necessary: omitting phi_D decouples debt from quality;
%   omitting psi_D decouples debt from speed; omitting eta_rev removes the pathway
%   by which under-staffed Code Review saturates human capacity.
% - Figure 4: plots of all four functions over their domain showing linear/saturation
%   regimes. Parameter default values annotated.

% ---- 3.7 ----
\subsection{Objective Functions and Cost Model}
\label{sec:model:objectives}
%
% Three Pareto objectives:
%
% V (maximize) = F_completed(T) / T
%   -- average features shipped per sprint over horizon T.
%
% Q (minimize) = B_escaped(T) / max(1, F_completed(T))
%   -- escaped bugs per feature delivered; quality proxy.
%
% C_F (minimize) = C_total / max(1, F_completed(T))
%   -- total cost per feature; cost-efficiency proxy.
%
% Cost components of C_total:
%   C_human  = (N * c_salary + c_PO) * (T / 52)      [annual salary, pro-rated]
%   C_agents = sum_k a_k * c_token_k * T              [per-sprint token costs]
%   C_bugs   = sum fixes: minor*C_fix_min + major*C_fix_maj
%   C_prod   = sum releases: B_escaped * C_prod_impact [fat-tailed; dominates]
%   C_total  = C_human + C_agents + C_bugs + C_prod
%
% - State the tri-objective problem: find A* in the feasible composition space
%   such that no configuration dominates on all three objectives simultaneously.
% - Objectives are treated symmetrically in hypervolume calculations;
%   V maximised, Q and C_F minimised.


% ============================================================
\section{Stability Analysis}
\label{sec:stability}
% ============================================================

% ---- 4.1 ----
\subsection{Overview and Purpose}
\label{sec:stability:overview}
%
% - Six stability conditions serve two purposes:
%   (1) Theoretical guarantees: prove model behaviour is bounded in limiting cases;
%       prevents degenerate configurations from inflating simulation runtimes.
%   (2) Feasibility pre-filter: conditions evaluated at prior means (expected-value
%       analysis) before any Monte Carlo run; infeasible configurations discarded.
% - Pre-filter is conservative (sufficient, not tight): some flagged-infeasible
%   configurations may be stable in practice due to favourable parameter draws.
%   Quantify false-positive rate in Section 7.
% - Conditions are checked sequentially (Steps 0-5 in src/feasibility.py);
%   first failure terminates the check.

% ---- 4.2 ----
\subsection{Theorem 1: Feature Backlog Stability}
\label{sec:stability:backlog}
%
% Statement: F_backlog(t) is stable (non-diverging) iff, at steady-state debt D*:
%   r_F * psi_D(D*) * a_3 + E[H_impl] >= lambda_F
%   i.e., expected implementation throughput >= feature arrival rate.
%
% Proof sketch:
%   - Model F_backlog as a discrete-time G/G/1 queue.
%   - Apply Loynes stability criterion; effective service rate is
%     r_F * psi_D(D*) * a_3 (debt-slowed), plus human residual coding H_impl.
%   - D* is the steady-state debt level from Theorem 4; substituting D* yields
%     the binding constraint on a_3.
%
% Practical implication: minimum Feature Implementation agents is a_3 >= lambda_F /
%   (r_F * psi_D(D*)). Violated frequently in configurations with large M and few
%   implementation agents.

% ---- 4.3 ----
\subsection{Theorem 2: Bug Queue Stability}
\label{sec:stability:bugs}
%
% Statement: bug queue B_known is stable iff:
%   (1) traffic intensity rho_bugs < 1: (net detectable bug rate) / (fix rate) < 1.
%   (2) regression feedback: p_regress * p_det_B < 1 (fixes don't self-sustain bugs).
%
% Proof sketch:
%   - Bugs arrive to B_known from latency buffer at rate lambda_det = lambda_inj * p_det_B.
%   - Fix throughput mu_fix = r_fix * psi_D(D*) * a_4.
%   - Regression injection adds p_regress * mu_fix bugs back per sprint;
%     net arrival = lambda_det + p_regress * mu_fix.
%   - Stability requires: (lambda_det + p_regress * mu_fix) / mu_fix < 1
%     => lambda_det / mu_fix < 1 - p_regress => rho_bugs < 1 - p_regress.
%
% Practical implication: most commonly triggered pre-filter condition, especially
%   in high-debt configurations with under-staffed Bug Fixing agents.

% ---- 4.4 ----
\subsection{Theorem 3: Human Standstill Condition}
\label{sec:stability:human}
%
% Statement: a configuration causes human standstill iff:
%   N * H_capacity <= N * H_comm(N) + N*M * H_monitoring_eff + H_review + H_escalation
%   i.e., monitoring + coordination + queue service exhausts all human capacity.
%
% Proof: from priority-order budget allocation (Step 7 of transition); if residual
%   capacity H_impl <= 0, humans perform no direct coding and no architecture
%   proposals are reviewed => Q_review grows unboundedly => no releases occur.
%
% Practical implication: provides explicit upper bound on M:
%   M_max = (H_capacity - H_comm(N) - H_review - H_escalation) / H_monitoring_eff
%   Configurations violating this are eliminated in Step 1 of pre-filter.
%   Interaction: high eta_esc (few Code Review agents) + high debt both increase
%   H_escalation, tightening M_max in a non-obvious way.

% ---- 4.5 ----
\subsection{Theorem 4: Debt Spiral Condition}
\label{sec:stability:debt}
%
% Statement: D(t) converges to finite D* in [0,1] iff:
%   delta_refactor * a_5 * r_refactor >
%   delta_impl * (r_F * psi_min * a_3 + E[H_impl]) + delta_fix * a_4 * r_fix_mean
%   i.e., refactoring removal rate > maximum debt injection rate.
%
% Proof:
%   - D(t+1) = D(t) + delta_impl*impl(t) - delta_refactor*refactor(t) + delta_fix*fixes(t)
%   - At D=1 (maximum): psi_D(1) = psi_min (minimum throughput); check if removal
%     still exceeds injection at this worst case.
%   - If condition holds, unique steady-state D* satisfies the fixed-point equation;
%     otherwise D diverges to 1 and throughput collapses.
%
% Practical implication: minimum Refactoring agent allocation scales with
%   Feature Implementation count; ratio a_5/a_3 must exceed a scenario-specific
%   threshold. Yields explicit formula for D*.

% ---- 4.6 ----
\subsection{Theorem 5: CI/CD Latency Queue Stability}
\label{sec:stability:latency}
%
% Statement: latency buffer L(t) remains bounded in expectation because tau_CI is
%   finite; bugs age out deterministically. However, VARIANCE of hidden bugs
%   E[|L|] = lambda_inj * tau_CI grows with tau_CI and bursty injection.
%
% Practical implication: motivates warm-up period T_warmup >= tau_CI_max before
%   metric collection begins. Also motivates DevOps agent allocation: more DevOps
%   agents (a_{11}) reduce tau_CI_base, shrinking E[|L|] and thus shortening the
%   tail of undetected bugs at release events.

% ---- 4.7 ----
\subsection{Composite Feasibility Theorem}
\label{sec:stability:composite}
%
% Statement: a configuration A is feasible iff it satisfies Theorems 1-5 simultaneously
%   at prior parameter means; violation of any single condition is sufficient for infeasibility.
%
% - Table 2: theorem summary -- name, condition expression, violated-by scenario,
%   pre-filter step index, frequency of violation across random samples (from Results).
%
% - Tightness discussion: prior-mean evaluation is conservative; Sobol' index
%   screening (Section 6.3) identifies which parameters most affect the probability
%   of borderline configurations being spuriously flagged infeasible.
%
% - The pre-filter reduces the effective search space; evolutionary search
%   operates only over the feasible polytope.


% ============================================================
\section{Simulation Methodology}
\label{sec:methodology}
% ============================================================

% ---- 5.1 ----
\subsection{Monte Carlo Protocol}
\label{sec:methodology:mc}
%
% - Per configuration: K independent replications, each with a distinct RNG stream.
% - Reproducibility: PCG-64 generator; stream seed = SeedSequence(master_seed,
%   spawn_key=(FNV-1a hash of config tuple, replication index k)).
%   Guarantees: (1) independent streams per (config, k) pair;
%               (2) identical result given same (config, K, seed).
%
% - Simulation horizon T sprints; warm-up T_warmup (>= tau_CI_max) sprints discarded
%   before metric collection.
%
% - Per replication: scenario parameters sampled once at sprint 0 from distributions
%   in Section 5.2 (lognormal, beta, uniform, gamma, or fixed); same parameters
%   held fixed throughout that replication's T sprints.
%
% - Aggregation: V, Q, C_F computed as means over (T - T_warmup) * K sprint-pairs;
%   95\% bootstrap confidence intervals (1000 bootstrap samples).
%
% - Convergence criterion for K: simulations run until CV of V (across K) < epsilon_CV;
%   report distribution of K achieved.
%
% - Table 3: simulation hyperparameters (T, T_warmup, K_max, epsilon_CV, RNG scheme).

% ---- 5.2 ----
\subsection{Parameter Distributions}
\label{sec:methodology:params}
%
% - ~74 scenario-level parameters drawn per replication from parameterised distributions:
%   * Throughput rates (r_F, r_fix, r_refactor, etc.): Log-normal (right-skewed;
%     most tasks fast, occasional long ones).
%   * Feature arrival rate lambda_F: Gamma (natural conjugate for Poisson arrival).
%   * Bug injection rate lambda_impl: Gamma (overdispersed positive rate).
%   * Detection rates (d_CR, d_SA, d_UT, d_IT, d_ST): Beta (bounded [0,1]).
%   * Fraction major bugs theta, p_regress, p_sec: Beta.
%   * Cost parameters: Log-normal (except c_salary, c_PO: fixed).
%   * CI/CD parameters (tau_CI_base, alpha_F_CI, alpha_test_CI): Log-normal or fixed.
%
% - Distribution justifications: reference empirical SE data where available
%   (defect injection rates from Basili et al., review detection rates from
%   Kemerer \& Paulk, CI pipeline latency from Hilton et al.); flag expert-elicited.
%
% - Three scenario parameterisations as shifts on the common baseline:
%   * Default: lambda_F = 2, balanced cost structure, tau_CI_base = 0.5.
%   * High-load: lambda_F = 4 (2x feature pressure), tau_CI_base = 0.75.
%   * Low-budget: c_token_k halved, c_salary unchanged (AI cost advantage).
%
% - Table 4: full parameter table with distribution family, baseline parameters,
%   scenario-specific overrides, and data source / justification.

% ---- 5.3 ----
\subsection{Search Strategies for Pareto-Optimal Configurations}
\label{sec:methodology:search}
%
% Two-phase search over the feasible composition space:
%
% Phase 1 -- Stratified random sampling:
%   - Sample N in [N_min, N_max], M in [1, M_max].
%   - Sample A as a random composition of N*M into 13 non-negative integers
%     using stars-and-bars (uniform over compositions), with repair to ensure a_3 >= 1.
%   - Apply feasibility pre-filter; discard infeasible configurations.
%   - Simulate feasible configurations; collect (V, Q, C_F) with bootstrap CIs.
%   - Output: initial Pareto approximation from n_random samples.
%
% Phase 2 -- Evolutionary search (initialised from Phase 1 non-dominated set):
%   - Mutation operator: move 1-3 agents between randomly chosen roles; repair
%     to maintain sum(A)=N*M and a_3>=1.
%   - Crossover: uniform crossover of two parent A vectors; repair constraints.
%   - Selection: NSGA-II crowding-distance tournament (3-objective hypervolume).
%   - Stopping: hypervolume improvement < epsilon_HV over last G generations,
%     or G_max generations reached.
%
% - Feasibility pre-filter applied at every candidate generation (both phases).
% - Hypervolume indicator tracks search progress; reference point = nadir of
%   initial random sample.
%
% - Figure 5: pseudocode / flowchart of the combined search procedure with
%   pre-filter integration.


% ============================================================
\section{Experimental Setup}
\label{sec:setup}
% ============================================================

% ---- 6.1 ----
\subsection{Scenarios}
\label{sec:setup:scenarios}
%
% Three scenarios representing distinct industrial contexts:
%
% Default (src/config/default_scenario.json):
%   - Steady-state product team; moderate feature demand lambda_F = 2.
%   - Balanced cost structure; tau_CI_base = 0.5 sprints.
%   - Representative of a mid-size agile product team.
%
% High-load (config/high_load.json):
%   - Startup-phase feature rush; lambda_F = 4 (2x default).
%   - Higher CI latency tau_CI_base = 0.75; larger codebase assumed.
%   - Tests whether optimal compositions for default generalise under pressure.
%
% Low-budget (config/low_budget.json):
%   - Resource-constrained team; AI token costs halved vs. default.
%   - Human salary unchanged; incentivises heavier AI-agent utilisation.
%   - Tests cost-efficiency frontier when API pricing is more favourable.
%
% - For each scenario provide the complete parameter vector (reference Table 4).
% - Justify scenario choices as covering a representative range of practical settings.

% ---- 6.2 ----
\subsection{Experimental Variables and Sweep Settings}
\label{sec:setup:sweep}
%
% - Ranges swept: N in {1,...,8}, M in {1,...,M_max(N)} where M_max from Theorem 3.
% - Composition space size: for total budget B = N*M agents across 13 roles,
%   number of compositions = C(B+12, 12) -- grows rapidly; motivates search.
% - Phase 1: n_random = 5000 configurations sampled per scenario.
% - Phase 2: pop_size = 100, G_max = 200 generations, epsilon_HV = 0.001.
% - K_max = 200 replications per configuration; T = 52 sprints; T_warmup = 5 sprints.
% - Hardware: [fill in after runs] CPU-hours, parallelised per configuration.
% - Software: Python 3.x, NumPy PCG-64, src/ simulation package.

% ---- 6.3 ----
\subsection{Evaluation Metrics and Comparisons}
\label{sec:setup:metrics}
%
% Primary metric: Pareto hypervolume dominated by discovered front per scenario.
%   - Computed with normalised objective values (nadir = reference point).
%   - Statistical test for scenario differences: Kruskal-Wallis + post-hoc.
%
% Secondary:
%   - Pairwise comparison of best-V, best-Q, best-C_F configurations.
%   - Role-composition profiles of Pareto-extreme solutions (Table 5).
%   - Fraction of configurations eliminated by each pre-filter step (Table 2).
%
% Ablation -- effect of degradation functions:
%   - Baseline: no-degradation model (phi_D = psi_D = 1, eta_rev = eta_esc = 1).
%   - Compare hypervolume and knee-point shift vs. full model (default scenario).
%   - Partial ablation: remove each of the four functions independently.
%
% Sensitivity analysis:
%   - Sobol' indices (first-order + total) for hypervolume w.r.t. all parameters.
%   - Report top-6 influential parameters as tornado chart.
%   - One-at-a-time +/- 20\% perturbation of top-6 to quantify Pareto-front shift.


% ============================================================
\section{Results}
\label{sec:results}
% ============================================================

% ---- 7.1 ----
\subsection{Pareto Frontiers by Scenario}
\label{sec:results:pareto}
%
% - Present tri-objective Pareto frontiers for all three scenarios:
%   3D scatter plot (V vs. Q vs. C_F) + three 2D projections (pairs of objectives).
%   Colour / marker distinguishes scenarios; CI-aware dominance used.
%
% - Report hypervolume values per scenario; statistical significance of differences.
%
% - Qualitative shape description:
%   * Knee point location (balanced V/Q/C_F).
%   * Steepness of quality-cost trade-off curve.
%   * Whether high-load front is dominated by default (throughput sacrifice).
%
% - Identify the practically interesting region (C_F within 10\% of min,
%   V >= 0.8 * V_max) and characterise it.
%
% - Figure 6: tri-objective Pareto scatter -- MAIN RESULT FIGURE.
%
% - Answer RQ1 in summary.

% ---- 7.2 ----
\subsection{Optimal Agent Compositions}
\label{sec:results:compositions}
%
% - Table 5: composition vector A for 4-5 Pareto-extreme solutions per scenario:
%   throughput-optimal, quality-optimal, cost-optimal, knee-point configuration.
%   Include objective values (V, Q, C_F) and confidence intervals.
%
% - Analysis -- which roles dominate quality-optimal configurations?
%   Hypothesis: a_6 (Code Review) + a_7 (Security Audit) + a_8-10 (Testing)
%   are over-allocated relative to throughput-optimal.
%
% - Role of Refactoring agents: do configurations with higher a_5/a_3 appear
%   more frequently near the interior (balanced) region of the Pareto front?
%
% - Fraction of configurations eliminated per pre-filter step; which condition
%   is most commonly triggered per scenario.
%
% - Span-of-control M: distribution of M across Pareto-optimal vs. dominated configs.

% ---- 7.3 ----
\subsection{Effect of Degradation Functions}
\label{sec:results:degradation}
%
% - Compare full model vs. no-degradation baseline (default scenario):
%   hypervolume ratio; shift in knee-point location.
%
% - Show that no-degradation model systematically overestimates V and
%   underestimates Q, confirming degradation functions are necessary.
%
% - Partial ablation: rank contribution of each degradation function to
%   hypervolume gap (phi_D, psi_D, eta_rev, eta_esc).
%
% - Figure 7: side-by-side Pareto front -- full model vs. no-degradation
%   baseline (default scenario).
%
% - Answer RQ2.

% ---- 7.4 ----
\subsection{Sensitivity Analysis}
\label{sec:results:sensitivity}
%
% - Sobol' indices: list top-6 influential parameters with first-order and
%   total-effect indices; discuss interaction effects.
%
% - Tornado chart: hypervolume sensitivity to +/- 20\% perturbation of top-6.
%   Parameters: likely candidates are lambda_F, lambda_impl, r_F, phi_max, psi_min,
%   p_regress.
%
% - Cross-scenario interpretation: high-load and low-budget scenarios as
%   perturbations along the most sensitive parameter directions.
%
% - Figure 8: tornado chart (or sensitivity scatter plots).
%
% - Answer RQ3.


% ============================================================
\section{Discussion}
\label{sec:discussion}
% ============================================================

% ---- 8.1 ----
\subsection{Interpretation of Findings}
\label{sec:discussion:interpretation}
%
% - Synthesise answers to RQ1, RQ2, RQ3 in plain language:
%   * RQ1: characterise the Pareto-optimal composition profile (e.g., heavy
%     Feature Implementation core + proportional Code Review + at least 1
%     Refactoring agent for balanced configurations).
%   * RQ2: quantify the performance gap introduced by degradation functions;
%     emphasise that ignoring technical debt dynamics leads to over-optimistic
%     throughput predictions and under-resourced testing allocations.
%   * RQ3: identify which aspects of optimal configurations are robust across
%     scenarios vs. which are scenario-specific.
%
% - Counterintuitive findings (if any): e.g., adding more agents does NOT
%   always improve throughput past M_max; highest-V configurations may not
%   be cost-optimal due to escalation overhead.
%
% - Connect to the formal theorems: optimal a_5/a_3 ratio predicted by
%   Theorem 4; M_max bound from Theorem 3 confirmed by empirical results.

% ---- 8.2 ----
\subsection{Practical Implications for AI-Human Team Design}
\label{sec:discussion:implications}
%
% - Translate Pareto-extreme composition profiles into actionable heuristics:
%   * Minimum Code Review + Testing allocation to avoid quality cliff in eta_rev.
%   * Refactoring-to-Feature-Implementation ratio implied by Theorem 4.
%   * Recommended span-of-control M range per scenario.
%
% - Span-of-control interaction: does higher M systematically shift the Pareto
%   front? Is there an optimal M* below M_max?
%
% - Tool/platform selection guidance: high-sensitivity parameters (Section 7.4)
%   identify the agent roles where capability improvements yield the largest
%   Pareto-front improvement; these are the highest-value targets for model upgrades.
%
% - Cost operationalisation: the per-agent per-sprint cost structure maps
%   directly to cloud LLM API pricing, making C_F directly actionable.
%   Compare cost at Pareto-optimal vs. a naive all-human baseline.

% ---- 8.3 ----
\subsection{Threats to Validity}
\label{sec:discussion:threats}
%
% INTERNAL VALIDITY
% - Monte Carlo convergence: report that CV < epsilon_CV met in all scenarios;
%   report actual K distribution.
% - Local optima in evolutionary search: compare evolutionary vs. random-search
%   hypervolume to bound sub-optimality gap.
% - Pre-filter conservatism: estimate false-positive rate by simulating a sample
%   of flagged-infeasible configurations.
%
% CONSTRUCT VALIDITY
% - Rate parameters drawn from assumed distributions (not validated against
%   specific industrial data); calibration strategy described in Section 5.2.
% - Degradation function forms (power law, piecewise linear) are assumed;
%   actual shapes may differ; ablation in Section 7.3 bounds the sensitivity.
% - Model does not capture inter-agent communication overhead or LLM context
%   limitations; these could be added as future extensions.
%
% EXTERNAL VALIDITY
% - Three scenarios may not cover all industrial contexts; N and M ranges finite.
% - Results may not generalise to very large teams (N > 10) where new bottlenecks
%   (e.g., inter-team coordination) emerge.
% - Single-project scope: multi-project portfolio settings introduce resource
%   sharing not modelled here.
%
% MODEL VALIDITY
% - 11-step sprint transition serialises steps that may overlap in practice.
% - Single-value tau_CI approximates what is in practice a distribution over
%   test execution times.
% - Product Owner modelled as singleton with unlimited capacity; unrealistic
%   at scale.
% - Formal stability theorems partially mitigate internal-validity concerns
%   by providing analytical guarantees in limiting cases.


% ============================================================
\section{Conclusion}
\label{sec:conclusion}
% ============================================================
%
% - Restate core contribution: first formal simulation framework for mixed
%   AI-human software team composition, grounded in a 14-dimensional sprint-state
%   model with proved stability guarantees and multi-objective optimisation.
%
% - Three main empirical findings (one sentence each, fill in after experiments).
%
% - Limitations not covered in threats:
%   * Static composition A; no adaptive re-allocation between sprints.
%   * No agent capability improvement over time (learning curves or model upgrades).
%   * Cost structure is static (no economy-of-scale effects).
%
% - Four concrete future-work directions:
%   1. Empirical calibration of rate parameters from industrial sprint logs and
%      LLM benchmark data (SWE-bench, HumanEval, CodeContests).
%   2. Multi-project portfolio extension: human developers shared across teams;
%      introduces cross-project capacity contention.
%   3. Online/adaptive composition policy: replace static A with a sprint-level
%      reallocation policy responding to observed S(t).
%   4. Agent capability growth: incorporate learning curves and model version
%      transitions (agent upgrades mid-project).


%\begin{acknowledgements}
% This work was supported by [funding agency] under grant [number].
%\end{acknowledgements}


% Authors must disclose all relationships or interests that
% could have direct or potential influence or impart bias on
% the work:
%
% \section*{Conflict of interest}
%
% The authors declare that they have no conflict of interest.


% ============================================================
% APPENDICES
% ============================================================
\appendix

\section{Proof Details}
\label{app:proofs}
%
% - Full proofs of Theorems 1-5 (proof sketches in Section 4 reference here).
% - Derivation of sprint transition equations from first principles.
% - Derivation of D* fixed-point (Theorem 4) and closed-form expression.
% - Proof that CI/CD latency buffer variance is bounded (Theorem 5).

\section{Parameter Tables}
\label{app:params}
%
% - Complete listing of all ~74 model parameters.
% - For each: symbol, description, distribution family, default parameters,
%   scenario-specific overrides (default / high-load / low-budget),
%   data source / justification.
% - Separate sub-table for team-composition control variables (N, M, A, R)
%   with sweep ranges used in experiments.


\bibliographystyle{spbasic}      % basic style, author-year citations
\bibliography{literature}        % name your BibTeX data base

\end{document}
